%============================================================
% FORMATTING: BASE LAYER (SHARED)
% Aturan yang 100% sama di Proposal, Progres, dan Buku PA.
% Aturan per-dokumen (geometry, spacing) ada di formats/format_*.tex
% yang di-input SETELAH file ini di masing-masing main_*.tex.
%============================================================

% ── Warna Institusi ──────────────────────────────────────────
\definecolor{pensbiru}{RGB}{0, 47, 167}    % Biru PENS (#002FA7)
\definecolor{penskuning}{RGB}{255, 193, 7} % Kuning/emas garis cover

% ── Nonaktifkan Hyphenation (pemenggalan kata) ──────────────
% Mencegah LaTeX memotong kata di akhir baris (mis. "Elek-tronika").
% tolerance=9999 + emergencystretch memberi ruang fleksibel agar
% LaTeX tidak perlu Overfull saat hyphenation dinonaktifkan.
\hyphenpenalty=10000
\exhyphenpenalty=10000
\tolerance=9999
\emergencystretch=3em

% ── Indentasi & Paragraf ─────────────────────────────────────
% \parindent TIDAK diset di sini - tiap dokumen punya nilai sendiri:
%   Buku PA  : 1.27cm (ada di formats/format_buku.tex)
%   Proposal : 0.75cm (ada di formats/format_proposal.tex)
%   Progres  : 0.75cm (ada di formats/format_progres.tex)
\setlength{\parskip}{0pt}

% ── Persamaan Matematika ─────────────────────────────────────
\setlength{\mathindent}{10mm}
\setlength{\abovedisplayskip}{12pt}
\setlength{\belowdisplayskip}{12pt}
\setlength{\abovedisplayshortskip}{12pt}
\setlength{\belowdisplayshortskip}{12pt}

% Nomor persamaan dicetak tebal-miring
\makeatletter
\renewcommand{\tagform@}[1]{\maketag@@@{\bfseries\itshape(#1)\@@italiccorr}}
\makeatother

% ── Format Judul Bab ─────────────────────────────────────────
\titleformat{\chapter}[display]
{\centering\fontsize{14}{17}\bfseries}
{BAB \thechapter}
{0pt}
{}
\titlespacing*{\chapter}{0pt}{-20pt}{18pt}

% ── Auto-selaraskan indentasi paragraf pertama dengan huruf judul (bukan nomor) ──
\newsavebox{\KotakLabelJudul}
\newdimen\LebarLabelJudul
\newcommand{\SelaraskanIndentAsli}[1]{%
  \sbox{\KotakLabelJudul}{#1}%
  \global\LebarLabelJudul=\wd\KotakLabelJudul
  \global\advance\LebarLabelJudul by 1em
  \xdef\ParindentAsliTersimpan{\the\parindent}%
  \global\parindent=\LebarLabelJudul
  \global\everypar{\parindent=\ParindentAsliTersimpan\relax\global\everypar{}}%
}

% ── Format Judul Subbab (Section) ────────────────────────────
\titleformat{\section}
{\fontsize{14}{17}\bfseries}
{\thesection}
{1em}
{}
[\SelaraskanIndentAsli{\thesection}]
\titlespacing*{\section}{0pt}{15pt}{0pt}

% ── Format Judul Sub-subbab (Subsection) ─────────────────────
\titleformat{\subsection}
{\fontsize{12}{14}\bfseries}
{\thesubsection}
{1em}
{}
[\SelaraskanIndentAsli{\thesubsection}]
\titlespacing*{\subsection}{0pt}{8pt}{3pt}

% ── Batas Penomoran dan Daftar Isi ───────────────────────────
% 0 = Chapter, 1 = Section, 2 = Subsection, 3 = Subsubsection
\setcounter{secnumdepth}{2}
\setcounter{tocdepth}{2}

% ── Format Judul Sub-sub-subbab (Subsubsection) ──────────────
% Tanpa nomor (secnumdepth=2) - heading ringan, spacing kecil.
% Label dikosongkan, sep 0pt, left=\parindent agar indent konsisten.
\titleformat{\subsubsection}
{\fontsize{12}{14}\bfseries}
{}
{0pt}
{}
\titlespacing*{\subsubsection}{\parindent}{3pt}{2pt}

% ── Caption Gambar ───────────────────────────────────────────
\captionsetup[figure]{
  labelfont     = {it},
  textfont      = {it},
  font          = {small},   % ~11pt dalam dokumen 12pt
  justification = centering,
  skip          = 6pt,
  name          = Gambar,
  labelsep      = space
}

% ── Caption Tabel ────────────────────────────────────────────
\captionsetup[table]{
  labelfont     = {it},
  textfont      = {it},
  font          = {small},
  justification = centering,
  skip          = 6pt,
  position      = above,
  name          = Tabel,
  labelsep      = space
}
\renewcommand{\tablename}{Tabel}

% ── Format Default Tabularray ────────────────────────────────
% Berlaku untuk semua lingkungan tblr dan longtblr
% row{1} styling TIDAK diset di sini agar Buku PA bisa pakai default sendiri.
% Proposal & Progres mengoverride via format_proposal.tex / format_progres.tex.
\SetTblrInner{
  hlines, vlines,                                             % Garis kotak otomatis
  rowhead = 1,                                               % Ulangi header jika tabel terpotong halaman
}

\SetTblrInner[longtblr]{
  hlines, vlines,
  rowhead = 1,
}

% ── Kustomisasi Khusus longtblr ──────────────────────────────
% 1. Menyamakan format caption longtblr dengan tabel biasa (Italic, 11pt, rata tengah)
\DefTblrTemplate{caption-sep}{default}{\space}   % Mengganti titik dua (:) menjadi spasi
\SetTblrStyle{caption-tag}{font=\small\itshape}  % "Tabel 3.1" miring & 11pt
\SetTblrStyle{caption-text}{font=\small\itshape} % Judul caption miring & 11pt
\SetTblrStyle{caption-halign}{c}                 % Memaksa caption rata tengah

% Mematikan teks "Continued on next page" dan caption ulangan di halaman selanjutnya
\DefTblrTemplate{contfoot-text}{default}{} % Hapus "Continued on next page"
\DefTblrTemplate{conthead-text}{default}{} % Hapus "(Continued)"
\DefTblrTemplate{capcont}{default}{}       % Hapus judul tabel ulangan
\DefTblrTemplate{contfoot}{default}{}      % Hapus ruang kosong di bawah
\DefTblrTemplate{conthead}{default}{}      % Hapus ruang kosong di atas

% ── Penomoran Gambar/Tabel/Persamaan per Bab ─────────────────
% Contoh: Gambar 3.1, Tabel 4.2, Persamaan (3.1)
\renewcommand{\thefigure}{\thechapter.\arabic{figure}}
\renewcommand{\thetable}{\thechapter.\arabic{table}}
\renewcommand{\theequation}{\thechapter.\arabic{equation}}

% ── Header & Footer ──────────────────────────────────────────
\pagestyle{fancy}
\fancyhf{}
\fancyfoot[C]{\thepage}
\renewcommand{\headrulewidth}{0pt}

\fancypagestyle{plain}{
  \fancyhf{}
  \fancyfoot[C]{\thepage}
  \renewcommand{\headrulewidth}{0pt}
}

% ── Lingkungan Kutipan Panjang (> 2 baris) ───────────────────
% Font 10pt, margin kanan-kiri masuk 10mm, ditulis miring
\newenvironment{kutipanpanjang}{%
  \begin{list}{}{%
      \setlength{\leftmargin}{10mm}
      \setlength{\rightmargin}{10mm}
      \setlength{\topsep}{6pt}
      \setlength{\itemsep}{0pt}
      \setlength{\parsep}{0pt}
    }
  \item\fontsize{10}{12}\selectfont\ignorespaces
  }{%
    \unskip
  \end{list}
}

% ── Hyperref ─────────────────────────────────────────────────
% Tautan hitam tanpa kotak berwarna
\hypersetup{
  colorlinks = true,
  linkcolor  = black,
  citecolor  = black,
  urlcolor   = black,
  pdftitle   = {\JudulPA},
  pdfauthor  = {\NamaMahasiswa}
}

% ── Nama Bahasa Indonesia ────────────────────────────────────
\renewcommand{\contentsname}{DAFTAR ISI}
\renewcommand{\listfigurename}{DAFTAR GAMBAR}
\renewcommand{\listtablename}{DAFTAR TABEL}
\renewcommand{\bibname}{DAFTAR PUSTAKA}
\renewcommand{\appendixname}{LAMPIRAN}
\renewcommand{\chaptername}{BAB}

% ── Judul Daftar Isi/Gambar/Tabel: 14pt, Bold, Center ────────
\renewcommand{\cfttoctitlefont}{\hfill\fontsize{14}{17}\bfseries\MakeUppercase}
\renewcommand{\cftaftertoctitle}{\hfill\null}

\renewcommand{\cftloftitlefont}{\hfill\fontsize{14}{17}\bfseries\MakeUppercase}
\renewcommand{\cftafterloftitle}{\hfill\null}

\renewcommand{\cftlottitlefont}{\hfill\fontsize{14}{17}\bfseries\MakeUppercase}
\renewcommand{\cftafterlottitle}{\hfill\null}

\setlength{\cftbeforetoctitleskip}{-20pt}
\setlength{\cftbeforeloftitleskip}{-20pt}
\setlength{\cftbeforelottitleskip}{-20pt}
\setlength{\cftaftertoctitleskip}{12pt}
\setlength{\cftafterloftitleskip}{12pt}
\setlength{\cftafterlottitleskip}{12pt}

% ── Prefix "Gambar" dan "Tabel" di LOF/LOT ───────────────────
\renewcommand{\cftfigpresnum}{Gambar }
\renewcommand{\cfttabpresnum}{Tabel }
\addtolength{\cftfignumwidth}{3.5em}
\addtolength{\cfttabnumwidth}{2.5em}

% ── Titik-titik di Daftar Isi ────────────────────────────────
\renewcommand{\cftdotsep}{0.5}

% ── Spasi Konsisten di Daftar Isi ────────────────────────────
\setlength{\cftbeforechapskip}{0pt}
\setlength{\cftbeforesecskip}{0pt}
\setlength{\cftbeforesubsecskip}{0pt}
\setlength{\cftbeforefigskip}{0pt}
\setlength{\cftbeforetabskip}{0pt}

% ── Font Level Bab: Bold ─────────────────────────────────────
\renewcommand{\cftchapfont}{\bfseries}
\renewcommand{\cftchappagefont}{\bfseries}
\renewcommand{\cftchapleader}{\bfseries\cftdotfill{\cftdotsep}}

% ── Font Level Section: Normal ───────────────────────────────
\renewcommand{\cftsecfont}{\normalfont}
\renewcommand{\cftsecpagefont}{\normalfont}

% ── Prefix "BAB" di Daftar Isi ───────────────────────────────
\renewcommand{\cftchappresnum}{BAB }
\addtolength{\cftchapnumwidth}{2.3em}

% ── Jarak (Spasi) Atas/Bawah Gambar dan Tabel ────────────────
\setlength{\intextsep}{12pt}
\setlength{\textfloatsep}{12pt}

\setlength{\bibsep}{3pt}

% ── Halaman Kosong ───────────────────────────────────────────
% Halaman kosong pada cetak bolak-balik diberi keterangan
\makeatletter
\newif\if@mainmatter
\@mainmatterfalse
\def\cleardoublepage{%
  \clearpage
  \if@twoside
  \ifodd\c@page\else
  \hbox{}
  \vspace*{\fill}
  \begin{center}
    \textit{Halaman ini sengaja dikosongkan}
  \end{center}
  \vspace{\fill}
  \if@mainmatter
  \thispagestyle{fancy}  % sudah masuk konten, tampilkan nomor
  \else
  \thispagestyle{empty}  % masih prelim sebelum roman, sembunyikan
  \fi
  \newpage
  \if@twocolumn\hbox{}\newpage\fi
  \fi
  \fi
}

% ── Paksa Halaman Ganjil (oneside & twoside) ─────────────────
% \cleartooddpage: lompat ke halaman ganjil berikutnya.
% Jika halaman saat ini sudah ganjil, tidak melakukan apa-apa.
% Jika genap, sisipkan halaman kosong bertuliskan keterangan
% lalu lanjutkan ke halaman ganjil berikutnya.
\newcommand{\cleartooddpage}{%
  \clearpage
  \ifodd\c@page\else
  \hbox{}
  \vspace*{\fill}
  \begin{center}
    \textit{Halaman ini sengaja dikosongkan}
  \end{center}
  \vspace{\fill}
  \thispagestyle{fancy}
  \newpage
  \fi
}
\makeatother
